% 官方第七条：按科技论文规范列出；正文引用处以 \cite 上标标注（GB/T 7714）
\begin{thebibliography}{99}
\bibitem{ref1} 陶文铨. 传热学[M]. 5版. 北京: 高等教育出版社, 2019.
\bibitem{ref2} 杨历, 陶斌斌. 多孔介质干燥过程传热传质研究[J]. 农业工程学报, 2005, 21(1): 27-31.
\bibitem{ref3} 湖北工业大学工程技术学院. 化工原理: 第八章 干燥[Z]. 武汉: 湖北工业大学工程技术学院.
\bibitem{ref4} 徐娓, 丁静, 赵义, 等. 多孔物料干燥中物料体积的收缩特性[J]. 华南理工大学学报(自然科学版), 2006, 34(8): 61-65.
\bibitem{ref5} 种翠娟, 朱文学, 刘云宏, 等. 胡萝卜薄层干燥动力学模型研究[J]. 食品科学, 2014, 35(9): 24.
\bibitem{ref6} 薛韩玲, 等. 基于湿度控制的大红袍花椒热风干燥动力学与品质研究[J]. 食品与发酵工业, 2023, 49(22): 149-155.
\bibitem{ref7} 任广跃, 朱乐雯, 段续, 等. 苹果丁冷冻-热风联合干燥体积收缩机制[J]. 农业工程学报, 2024, 40(2): 63-71.
\bibitem{ref8} 顾祥红. 无限长圆柱非稳态导热与集总参数法应用[J]. 沈阳理工大学学报, 2009(2).
\bibitem{ref9} GASPARIN S, BERGER J, DUTYKH D, et al. On the solution of coupled heat and moisture transport in porous material[J]. Transport in Porous Media, 2018, 121: 665-702.
\bibitem{ref10} HUANG F, SHEN J. On a new class of BDF and IMEX schemes for parabolic type equations[J]. SIAM Journal on Numerical Analysis, 2024, 62(4): 1609-1637.
\bibitem{ref11} ISLAM M, HYE A, MAMUN A. Nonlinear effects on the convergence of Picard and Newton iteration methods in the numerical solution of one-dimensional variably saturated–unsaturated flow problems[J]. Hydrology, 2017, 4(4): 50.
\bibitem{ref12} ADROVER A, BRASIELLO A, PONSO G. A moving boundary model for food isothermal drying and shrinkage: general setting[J]. Journal of Food Engineering, 2019, 244: 178-191.
\bibitem{ref13} ADROVER A, BRASIELLO A, PONSO G. A moving boundary model for food isothermal drying and shrinkage: a shortcut numerical method for estimating the shrinkage factor[J]. Journal of Food Engineering, 2019, 244: 212-219.
\bibitem{ref14} ADROVER A, BRASIELLO A. A moving boundary model for food isothermal drying and shrinkage: one-dimensional versus two-dimensional approaches[J]. Journal of Food Process Engineering, 2019, 42(6): e13178.
\bibitem{ref15} SADOUN N, SI-AHMED E K, COLINET P, et al. On the boundary immobilization and variable space grid methods for transient heat conduction problems with phase change: discussion and refinement[J]. Comptes Rendus Mécanique, 2012, 340(7): 501-511.
\bibitem{ref16} CRANK J. Two methods for the numerical solution of moving-boundary problems in diffusion and heat flow[J]. The Quarterly Journal of Mechanics and Applied Mathematics, 1957, 10(2): 220-231.
\bibitem{ref17} MURRAY W D, LANDIS F. Numerical and machine solutions of transient heat-conduction problems involving melting or freezing: Part I—method of analysis[J]. Journal of Heat Transfer, 1959, 81(2): 106-114.
\bibitem{ref18} GIBSON R E. A heat conduction problem involving a specified moving boundary[J]. Quarterly of Applied Mathematics, 1959, 16(4).
\bibitem{ref19} 胥文博. 热传导方程两类参数的统计反演方法[J]. 应用数学进展, 2025, 14(3): 203-215.
\bibitem{ref20} 王清艳. 扩散方程反问题的正则化方法[J]. 理论数学, 2024, 14(4): 98-106.
\bibitem{ref21} Optimal Adaptive Implicit Time Stepping[EB/OL]. arXiv:2506.18809, 2025.
\end{thebibliography}
